% ===========================================================================
%  comune/stile-tikz.tex — il vocabolario grafico delle figure
% ===========================================================================
%
%  Ogni figura del progetto disegna con QUESTI stili e non con parametri
%  propri. La ragione e' leggibilita': se in una figura il riquadro di un
%  modulo e' azzurro e in quella dopo e' dello stesso azzurro un archivio su
%  disco, chi legge impara due volte e si fida zero. Qui un colore significa
%  una cosa sola in tutto il documento:
%
%     cielo     il MOTORE di calcolo (Python, core/)
%     menta     l'INTERFACCIA (Flask, dashboard/)
%     pesca     il SITO (Next.js) e tutto cio' che sta in rete
%     lavanda   i dati PERSISTENTI (disco, Postgres)
%     crema     cio' che viene da FUORI (broker, Stripe, provider dati)
%     bianco    strati astratti (porte, contratti, protocolli)
%
%  Questo file e' incluso sia dai documenti sia dalle figure compilate da sole
%  (figure/tikz/*.tex con la classe standalone), percio' non deve dipendere da
%  niente che esista solo dentro un documento.
% ===========================================================================
\RequirePackage{tikz}
\usetikzlibrary{
  arrows.meta, positioning, fit, backgrounds, calc, shapes.geometric,
  shapes.multipart, shapes.symbols, decorations.pathreplacing,
  decorations.pathmorphing, patterns, shadows.blur, matrix, chains, scopes
}

\tikzset{
  % --- impostazioni generali ----------------------------------------------
  every picture/.style = {
    line width = 0.7pt,
    font = \sffamily\small,
  },
  %
  % --- RIQUADRI: un modulo, un componente, una responsabilita' -------------
  riquadro/.style = {
    rectangle, rounded corners = 3pt, draw,
    minimum height = 9.5mm, minimum width = 30mm,
    inner xsep = 6pt, inner ysep = 5pt,
    align = center, font = \sffamily\small,
    text = lyraInchiostro,
  },
  motore/.style    = {riquadro, fill = lyraCielo,   draw = lyraCieloBordo},
  interfaccia/.style = {riquadro, fill = lyraMenta,  draw = lyraMentaBordo},
  sito/.style      = {riquadro, fill = lyraPesca,   draw = lyraPescaBordo},
  persistenza/.style = {riquadro, fill = lyraLavanda, draw = lyraLavandaBordo},
  esterno/.style   = {riquadro, fill = lyraCrema,   draw = lyraCremaBordo},
  astratto/.style  = {riquadro, fill = white, draw = lyraGrigio, dashed},
  %
  % Il riquadro piccolo: per gli elenchi fitti dove il nome basta.
  minuto/.style = {riquadro, minimum height = 7mm, minimum width = 22mm,
                   font = \sffamily\scriptsize, inner xsep = 4pt},
  %
  % --- FORME PARTICOLARI ---------------------------------------------------
  archivio/.style = {
    cylinder, shape border rotate = 90, aspect = 0.22, draw = lyraLavandaBordo,
    fill = lyraLavanda, minimum height = 13mm, minimum width = 22mm,
    inner xsep = 0pt, inner ysep = 4pt, align = center,
    font = \sffamily\scriptsize, text = lyraInchiostro,
  },
  decisione/.style = {
    diamond, draw = lyraInchiostro!60, fill = white, aspect = 1.9,
    inner xsep = 2pt, inner ysep = 0pt, align = center,
    font = \sffamily\scriptsize, text = lyraInchiostro,
  },
  persona/.style = {
    rectangle, rounded corners = 10pt, draw = lyraInchiostro!70,
    fill = white, minimum height = 9mm, minimum width = 24mm,
    align = center, font = \sffamily\small\itshape,
  },
  %
  % --- STRATI: il contenitore che raggruppa -------------------------------
  strato/.style = {
    rectangle, rounded corners = 5pt, draw = lyraGrigioChiaro,
    fill = lyraPannello!55!white, inner sep = 9pt,
  },
  stratoTitolo/.style = {
    font = \sffamily\scriptsize\bfseries, text = lyraInchiostro!65,
    anchor = north west,
  },
  %
  % --- FRECCE --------------------------------------------------------------
  %  Tre tipi, e la differenza va vista senza leggere la legenda:
  %    flusso      pieno    -> i dati passano di qui
  %    dipendenza  tratteg. -> questo conosce quello (ma nulla scorre)
  %    ritorno     grigio   -> la risposta che torna indietro
  flusso/.style = {
    -{Stealth[length = 2.4mm, width = 1.8mm]},
    draw = lyraInchiostro!75, line width = 0.8pt,
  },
  dipendenza/.style = {
    -{Stealth[length = 2.2mm, width = 1.6mm]}, dashed, dash pattern = on 2.4pt off 2pt,
    draw = lyraInchiostro!55, line width = 0.6pt,
  },
  ritorno/.style = {
    -{Stealth[length = 2.2mm, width = 1.6mm]},
    draw = lyraGrigio, line width = 0.6pt,
  },
  accento/.style = {
    -{Stealth[length = 2.6mm, width = 2mm]},
    draw = lyraAccento, line width = 1pt,
  },
  bloccato/.style = {
    -{Bar[width = 2.4mm]}, draw = lyraNegativo, line width = 0.9pt,
  },
  %
  % --- ETICHETTE SULLE FRECCE ---------------------------------------------
  etichetta/.style = {
    font = \sffamily\scriptsize, text = lyraInchiostro!70,
    fill = white, inner xsep = 2.5pt, inner ysep = 1pt,
  },
  etichettaAccento/.style = {etichetta, text = lyraAccentoScuro},
  %
  % --- NOTE A MARGINE ------------------------------------------------------
  nota/.style = {
    font = \sffamily\scriptsize\itshape, text = lyraInchiostro!60,
    align = left, inner sep = 2pt,
  },
  notaAccento/.style = {nota, text = lyraAccentoScuro, font = \sffamily\scriptsize},
  %
  % --- SEGNAPOSTO NUMERICI (1) (2) (3) sulle figure a passi ---------------
  passo/.style = {
    circle, draw = lyraAccento, fill = white, text = lyraAccentoScuro,
    font = \sffamily\scriptsize\bfseries, inner sep = 1.2pt,
    minimum size = 4.6mm,
  },
}

% --- la cornice con il titolo, usata dalle figure a strati -------------------
%  #1 opzioni  #2 nodi da contenere  #3 titolo
%  Il titolo sta SOPRA la cornice, non dentro: dentro finisce addosso al primo
%  nodo ogni volta che il nodo è vicino al bordo, e lo si scopre solo guardando
%  il PDF. Il nome del nodo è contato, così due strati nella stessa figura non
%  si sovrascrivono a vicenda.
\newcounter{lyrastratoN}
\newcommand{\lyrastrato}[3][]{%
  \stepcounter{lyrastratoN}%
  \begin{scope}[on background layer]
    \node[strato, fit = #2, #1] (cornice\thelyrastratoN) {};
    \node[stratoTitolo, anchor = south west]
      at ([shift={(1pt,1.5pt)}]cornice\thelyrastratoN.north west) {#3};
  \end{scope}}

% --- legenda compatta, da mettere sotto una figura ---------------------------
%  Si usa solo quando la figura introduce un colore nuovo: ripeterla sotto ogni
%  figura la fa diventare rumore che nessuno legge piu'.
\newcommand{\lyralegenda}[1]{%
  \begin{tikzpicture}[baseline]
    \matrix[matrix of nodes, nodes = {anchor = west, font = \sffamily\scriptsize},
            column sep = 3mm, row sep = 0.6mm] {#1\\};
  \end{tikzpicture}}

% ---------------------------------------------------------------------------
%  DUE TRAPPOLE, scritte qui perché costano mezz'ora ciascuna e l'errore che
%  producono indica sempre il punto sbagliato.
%
%  1. `\\` dentro un gruppo del contenuto di un nodo.
%        {testo\\{\scriptsize una riga\\e un'altra}}   <-- SI ROMPE
%        {testo\\{\scriptsize una riga}\\{\scriptsize e un'altra}}  <-- giusto
%     TikZ raccoglie le righe guardando i `\\` al primo livello: quello dentro
%     le graffe è il `\\` di LaTeX, che in un box orizzontale non ha senso.
%     L'errore che ne esce è «\tikzscope@linewidth non definito», segnalato sul
%     nodo DOPO quello colpevole.
%
%  2. Una tcolorbox dentro un nodo. Vale per \cod, \pulsante, \campo: sono
%     tcbox, cioè figure TikZ annidate, e lasciano la figura esterna in uno
%     stato incoerente. Stesso errore, stessa confusione. Nelle figure si usa
%     \codT (qui sotto).
% ---------------------------------------------------------------------------

% --- il codice DENTRO una figura --------------------------------------------
%  Nel testo un simbolo si scrive \cod{...}, che è un \tcbox: una tcolorbox,
%  cioè una tikzpicture. Infilarla dentro il contenuto di un nodo TikZ rompe lo
%  stato della figura che la contiene, e l'errore salta fuori sul nodo
%  SUCCESSIVO (`\tikzscope@linewidth` non definito) — mezz'ora buttata a
%  cercarlo nel posto sbagliato. Dentro le figure si usa questo, che è solo un
%  cambio di carattere e non apre niente.
%  \upshape serve davvero: dentro un nodo `nota`, che è in corsivo, un
%  monospazio corsivo si legge male e sembra un'altra cosa.
\newcommand{\codT}[1]{{\ttfamily\upshape\small\color{lyraInchiostro!85}#1}}

\newcommand{\lyravoce}[2]{%
  \tikz[baseline=-0.55ex]\node[#1, minimum width = 4.2mm, minimum height = 2.6mm,
        inner sep = 0pt, rounded corners = 1.4pt]{};~#2}
